\usepackage[scr=rsfs,cal=euler]{mathalfa}
\multlinegap0pt

\hyphenation{begin decom-po-si-tion sequence}
\newenvironment{enumeratea}
{\bgroup\def\theenumi{\alph{enumi}}\begin{enumerate}}
{\end{enumerate}\egroup}

\newcommand{\ba}{{\boldsymbol a}}
\newcommand{\bb}{{\boldsymbol b}}
\newcommand{\bd}{{\boldsymbol d}}
\newcommand{\bu}{{\boldsymbol u}}
\newcommand{\bx}{{\boldsymbol x}}
\newcommand{\bmy}{{\boldsymbol y}}

\newcommand{\bD}{{\boldsymbol D}}
\newcommand{\bL}{{\boldsymbol L}}
\newcommand{\bP}{{\boldsymbol P}}
\newcommand{\bV}{{\boldsymbol V}}
\newcommand{\bX}{{\boldsymbol X}}
\newcommand{\bZ}{{\boldsymbol Z}}

\newcommand{\Psfrac}[2]{(\sfrac{#1}{#2})}
\newcommand{\psfrac}[2]{\sfrac{(#1)}{#2}}
\newcommand{\bigO}[1]{O(#1)}

\usepackage{mathtools}
\usepackage{bbm}

\newcommand{\NN}{\mathbb{N}}
\newcommand{\ZZ}{\mathbb{Z}}
\newcommand{\RR}{\mathbb{R}}
\newcommand{\CC}{\mathbb{C}}
\newcommand{\PP}{\mathbb{P}}
\newcommand{\EE}{\mathbb{E}}

\newcommand{\boldalpha}{\boldsymbol{\alpha}}
\newcommand{\bolddelta}{\boldsymbol{\delta}}
\newcommand{\boldrho}{\boldsymbol{\rho}}
\newcommand{\dee}{\,\mathrm{d}}
\newcommand{\one}{\mathbbm{1}}
\newcommand{\BC}{\mathcal{B}}
\newcommand{\EC}{\mathcal{E}}
\newcommand{\RS}{\mathscr{R}}
\newcommand{\prob}[1]{\PP[ #1 ]}
\newcommand{\bigoo}[1]{O\!\left(#1 \right)}
\newcommand{\abs}[1]{\left\lvert #1 \right\rvert}
\newcommand{\floor}[1]{\left\lfloor #1 \right\rfloor}
\newcommand{\riemint}[4]{\int_{#1}^{#2}\! #3 \, \mathrm{d}#4}
\newcommand{\maxnorm}[1]{\left\| #1 \right\|_{\infty}}

\renewcommand{\Re}{\textup{Re}}
\renewcommand{\Im}{\textup{Im}}

\newcommand{\dtv}{\mathrm{d}_{\mathrm{TV}}}
\newcommand{\dw}{\mathrm{d}_{\mathrm W}}

\let\epsilon\varepsilon

\newcommand{\pextra}{P_{\mathtt{extra}}}
\newcommand{\parratia}{P_{\mathtt{Arratia}}}

\newcommand{\simplex}{\Delta^{k-1}}

\newcommand{\bg}{\big}
\newcommand{\bgg}{\Big}
\newcommand{\bggg}{\bigg}

\newcommand{\Rnt}{R_{\mathrm{NT}}}
\newcommand{\Rprob}{R_{\mathrm{PD}}}
\newcommand{\ECnt}{\EC_{\mathrm{NT}}}
\newcommand{\ECprob}{\EC_{\mathrm{PD}}}

\newcommand{\Bnt}{\mathcal{B}_{\mathrm{NT}}}
\newcommand{\Bprob}{\mathcal{B}_{\mathrm{PD}}}

\DeclareMathOperator{\vard}{Var}
\DeclareMathOperator{\Dir}{Dir}

\theoremstyle{plain}
\newtheorem{cor}{Corollary}[section]
\newtheorem{lem}[cor]{Lemma}
\newtheorem{prop}[cor]{Proposition}
\newtheorem{conj}[cor]{Conjecture}
\newtheorem{theo}{Theorem}

\theoremstyle{definition}
\newtheorem{example}{Example}
\newtheorem{deff}[cor]{Definition}
\newtheorem{rem}{Remark}[section]
\newtheorem*{rems}{Remark}
\newtheorem*{remss}{Remarks}

\makeatletter
\ams@newcommand{\multiint}[1]{\DOTSI\protect\MultiIntegral{#1}}
\renewcommand{\MultiIntegral}[1]{%
\edef\ints@c{\noexpand\intop
\ifnum#1=\z@\noexpand\intdots@\else\noexpand\intkern@\fi
\replicate{#1-2}{\noexpand\intop\noexpand\intkern@}%
\noexpand\intop
\noexpand\ilimits@
}%
\futurelet\@let@token\ints@a
}
\makeatother
\ExplSyntaxOn
\cs_new:Npn \replicate #1 #2 { \prg_replicate:nn { #1 } { #2 } }
\ExplSyntaxOff
